\documentclass[letterpaper,10pt,conference]{ieeeconf}

\IEEEoverridecommandlockouts
\overrideIEEEmargins
\pdfminorversion=4

\usepackage{amsmath,amssymb}
\usepackage{booktabs}
\usepackage{graphicx}
\usepackage{microtype}
\usepackage{url}

\newcommand{\method}{BQ-Splat}
\newcommand{\R}{\mathbb{R}}

\title{BQ-Splat: Bayesian Quadrature for Post-hoc Uncertainty\\in Gaussian Splat Maps}

% RA-L uses double-anonymous review for initial submissions.
\author{Anonymous Authors}

\begin{document}

\maketitle
\thispagestyle{empty}
\pagestyle{empty}

\begin{abstract}

\end{abstract}

\begin{center}
\small\textbf{Index Terms---}Computer Vision for Automation; Mapping;
Perception and Sensing; Gaussian splatting; uncertainty quantification.
\end{center}

\section{Introduction}


\section{Background}

\subsection{Rendering with Gaussian splats}

Neural radiance fields represent density and view-dependent colour as
continuous functions and render pixels by integrating radiance along camera
rays \cite{mildenhall2020}. 3DGS instead represents the scene explicitly as
anisotropic Gaussian primitives \cite{kerbl2023}. Each primitive $i$ has a
centre $\mu_i$, covariance $\Sigma_i$, opacity parameter, and a
view-dependent colour $c_i(d)$, commonly represented by spherical harmonics.

For a pixel $q$, the visible Gaussians are projected to the image plane and
sorted from front to back. If $\alpha_{q,i}$ is the projected opacity of the
$i$th primitive and

\begin{equation}
    T_{q,i}=\prod_{j<i}(1-\alpha_{q,j})
\end{equation}

is its accumulated transmittance, the rendered colour is

\begin{equation}
    C(q)=\sum_{i=1}^{m} T_{q,i}\alpha_{q,i}c_i(d_q).
    \label{eq:alpha-compositing}
\end{equation}

Equation~\eqref{eq:alpha-compositing} is a weighted sum, but it is also a
quadrature rule. It evaluates the continuous volume-rendering functional

\begin{equation}
    \mathcal{L}_q[c]=\int a_q(\xi)c(\xi,d_q)\,d\xi,
    \label{eq:rendering-functional}
\end{equation}

where $a_q(\xi)$ combines transmittance, opacity or density, and the projected
footprint of the ray. Under the piecewise-constant colour and density model
used in volume rendering, the integral reduces exactly to the ordered
alpha-compositing sum \cite{max1995}. Thus compositing is not merely analogous
to quadrature: it is a particular quadrature rule for the rendering integral,
with nodes supplied by the primitives and weights supplied by the renderer.

\subsection{From quadrature rules to uncertainty}

Consider a general integral

\begin{equation}
    I[f]=\int_{\Omega} f(x)p(x)\,dx,
    \label{eq:integral}
\end{equation}

where $f$ is known only at locations $X=\{x_1,\ldots,x_m\}$. A quadrature
rule approximates this integral by

\begin{equation}
    I[f]\approx \widehat I[f]=\sum_{i=1}^{m}w_i f(x_i).
    \label{eq:quadrature}
\end{equation}

The nodes $x_i$ determine where the function is evaluated and the weights
$w_i$ determine how those values contribute. Newton--Cotes, Gaussian
quadrature, Monte Carlo integration, and alpha compositing all take this form;
they differ in how they choose nodes and weights. A conventional quadrature
rule returns an estimate but does not, by itself, express what could have
happened between the samples.

Probabilistic numerics treats that missing information as an inference
problem. In Bayesian quadrature \cite{ohagan1991,briol2019}, the unknown
integrand is assigned a Gaussian-process prior,

\begin{equation}
    f\sim\mathcal{GP}(0,k),
\end{equation}

where $k$ encodes assumptions about the variation of $f$. Conditioning on
$\mathbf{f}=[f(x_1),\ldots,f(x_m)]^\top$ produces a posterior over functions.
Since integration is linear, the induced posterior over the integral is
Gaussian:

\begin{equation}
    I[f]\mid \mathbf f,X \sim
    \mathcal N\!\left(\mathbf z^\top K^{-1}\mathbf f,
    z_0-\mathbf z^\top K^{-1}\mathbf z\right),
    \label{eq:bq-posterior}
\end{equation}

with

\begin{align}
    K_{ij} &= k(x_i,x_j), \\
    z_i &= \int_{\Omega}k(x,x_i)p(x)\,dx, \\
    z_0 &= \int_{\Omega}\!\int_{\Omega}
    k(x,x')p(x)p(x')\,dx\,dx'.
\end{align}

The posterior mean is itself a quadrature rule, with weights
$\mathbf w=K^{-1}\mathbf z$. The posterior variance

\begin{equation}
    u(\Omega)=z_0-\mathbf z^\top K^{-1}\mathbf z
    \label{eq:variance}
\end{equation}

depends on the sample locations, integration measure, and kernel, but not on
the observed function values. It is small when the samples constrain all
functions allowed by the prior over the relevant domain. It is large when
substantially different functions can pass through the same observations yet
produce different integrals. This is the sense in which Bayesian quadrature
provides epistemic uncertainty about a numerical computation.

The same quantity has a deterministic interpretation. For functions in the
reproducing-kernel Hilbert space associated with $k$,

\begin{equation}
    |I[f]-\widehat I[f]|
    \leq \lVert f\rVert_k\sqrt{u(\Omega)}.
    \label{eq:error-bound}
\end{equation}

The posterior standard deviation is therefore also the worst-case integration
error over the unit ball of that function space.

\section{Related Work}

\subsection{Uncertainty in radiance fields and Gaussian splatting}

Uncertainty methods for radiance fields and 3DGS broadly operate on model
parameters, rendered predictions, or observation coverage. Stochastic
Gaussian Splatting learns distributions over Gaussian parameters using
variational inference \cite{savant2024}. Bayes' Rays applies a Laplace
approximation to a deformation field \cite{goli2024}. FisherRF uses Fisher
information for uncertainty estimation and active view selection
\cite{jiang2024}, while PUP 3D-GS uses second-order sensitivity for pruning
\cite{hanson2025}. GAVIS records directional visibility using spherical
harmonics for active mapping \cite{xue2026}. Mip-Splatting addresses a related
but distinct source of error: aliasing under changes in image sampling rate
\cite{yu2024}.

Galappaththige et al.\ recently proposed a post-hoc method for predictive
photometric uncertainty in 3DGS \cite{galappaththige2026predictive}. Their
method freezes a trained map and learns an additional per-primitive
uncertainty channel from training-view reconstruction residuals. These
channels are represented with spherical harmonics and rendered with the same
alpha-compositing weights as colour. A Gaussian prior, implemented through
$L^2$ regularisation, drives unsupported directions towards a prescribed
maximum uncertainty.


\subsection{Bayesian quadrature for representation analysis}

Bayesian quadrature is usually used when function evaluations are expensive
and uncertainty in an integral matters \cite{ohagan1991,briol2019}. Our use is
slightly different. The function values stored by a trained 3DGS map are
already available; the question is whether their locations provide adequate
support for a new spatial or directional query. The posterior variance is
therefore used as a geometric property of the representation. This turns a
tool for uncertain numerical integration into a post-hoc diagnostic for
finite scene maps.

\section{Bayesian Quadrature for Gaussian Splat Maps}

\subsection{Spatial support}

Let $\Omega_q\subset\R^3$ be a small region centred on a query point $q$, and
let $X_q=\{x_1,\ldots,x_m\}$ be the centres of nearby splats. We interpret the
splat centres as sample locations of an unknown local scene function. With a
uniform integration density over $\Omega_q$, Eq.~\eqref{eq:variance} gives a
scalar measure of spatial support.

The prior term $z_0$ is the uncertainty in the local integral before any
splats are observed. The term $\mathbf z^\top K^{-1}\mathbf z$ is the
uncertainty removed by their locations. Nearby splats distributed across the
region produce a large reduction; few, distant, or clustered splats leave a
larger variance. The measure consequently captures sample arrangement rather
than only sample count.

We use an axis-aligned box for $\Omega_q$, retrieve neighbours with a $k$-d
tree, and cap each solve at 400 splats. For separable kernels and box domains,
the required moments factor into one-dimensional integrals. Since spatial
uncertainty uses only the primitive locations, it can be computed from a
trained map without colours, residuals, or additional optimisation.

\subsection{Directional support}

View-dependent appearance requires support in both position and direction. We
represent an observation by $(x_i,d_i)$, where $d_i\in\mathbb S^2$ is a camera
direction from which splat $i$ was observed, and use the product kernel

\begin{equation}
    k\big((x,d),(x',d')\big)
    =k_{\mathrm{pos}}(x,x')k_{\mathrm{dir}}(d,d').
    \label{eq:product-kernel}
\end{equation}

Our directional kernel is

\begin{equation}
    k_{\mathrm{dir}}(d,d')
    =\exp\!\left(\kappa(d^\top d'-1)\right),
    \label{eq:direction-kernel}
\end{equation}

where $\kappa$ controls the angular length scale. Small $\kappa$ transfers
information across a wide range of directions; large $\kappa$ represents
stronger view dependence.

For query direction $d_*$, the kernel mean is

\begin{equation}
    z_i=\left[\int_{\Omega_q}k_{\mathrm{pos}}(x,x_i)p(x)\,dx\right]
    k_{\mathrm{dir}}(d_i,d_*).
    \label{eq:directional-mean}
\end{equation}

The resulting variance is low only when the query is supported in both
position and direction. Position-only and joint position--direction
uncertainty are therefore the same construction applied to different input
spaces. We recover camera--splat associations from the training cameras using
frustum intersection and an angular depth test, then attach the corresponding
viewing direction to each visible splat.

\subsection{The rendering-aware kernel}
\label{sec:rendering-aware}

Uniform box quadrature asks whether splats support a local spatial average.
Rendering asks a more specific question: whether they support the weighted
radiance integral in Eq.~\eqref{eq:rendering-functional}. Let

\begin{equation}
    a_q(\xi)=T_q(\xi)\,\sigma(\xi)\,G_q(\xi)
    \label{eq:render-weight}
\end{equation}

denote the query-dependent rendering weight, combining accumulated
transmittance, density or opacity, and the projected footprint or visibility
term. There are two equivalent ways to incorporate this structure.

First, place a GP prior on radiance,
$c\sim\mathcal{GP}(0,k_{\mathrm{base}})$, and apply the linear rendering
functional $\mathcal L_q$ to that process. The Bayesian-quadrature moments are

\begin{align}
    z_{q,i} &= \int a_q(\xi)
    k_{\mathrm{base}}(\xi,x_i)\,d\xi, \\
    z_{q,0} &= \int\!\int a_q(\xi)a_q(\xi')
    k_{\mathrm{base}}(\xi,\xi')\,d\xi\,d\xi'.
    \label{eq:render-moments}
\end{align}

The observation Gram matrix remains
$K_{ij}=k_{\mathrm{base}}(x_i,x_j)$ because the observations are colours
$c(x_i)$.

Second, define the rendering integrand $g_q(\xi)=a_q(\xi)c(\xi)$. It is also a
Gaussian process, with the query-specific kernel

\begin{equation}
    k_q(\xi,\xi')=
    a_q(\xi)k_{\mathrm{base}}(\xi,\xi')a_q(\xi').
    \label{eq:render-kernel}
\end{equation}

Integrating $g_q$ against the ordinary volume measure gives exactly the same
$z_{q,0}$ as Eq.~\eqref{eq:render-moments}. Equation~\eqref{eq:render-kernel}
is the rendering-aware kernel structure: the renderer gates covariance at
both arguments. Points have little influence when they are occluded, have
negligible projected opacity, or lie outside the pixel footprint. This
equivalence also clarifies an implementation detail: when conditioning on
stored radiance values, use the base Gram matrix and rendering-weighted
moments; when conditioning on values of the weighted integrand, use the full
query-specific Gram matrix.

We approximate $a_q$ by an unnormalised Gaussian whose mass, mean, and
covariance match the contributing splats. For an RBF base kernel, all moments
then follow in closed form from Gaussian-product identities. We consider
three sources of rendering weights: mean local opacity, depth-ordered ray
transmittance, and the anisotropic per-pixel alpha values produced by
\texttt{gsplat}. In the last two cases, the total mass is
$\sum_iT_{q,i}\alpha_{q,i}=1-T_{q,\mathrm{final}}$, matching the standard
alpha-compositing recursion.

The directional extension follows by setting

\begin{equation}
    k_{\mathrm{base}}\big((x,d),(x',d')\big)
    =k_{\mathrm{pos}}(x,x')k_{\mathrm{dir}}(d,d').
\end{equation}

The rendering weight remains spatial: it describes where the ray receives
mass. Directional support enters through the base covariance. Consequently,

\begin{equation}
    z_{q,i}=\left[\int a_q(\xi)
    k_{\mathrm{pos}}(\xi,x_i)\,d\xi\right]
    k_{\mathrm{dir}}(d_i,d_*).
    \label{eq:render-directional-mean}
\end{equation}

This separates two distinct reasons for uncertainty: weak contribution along
the ray and poor coverage of the requested viewing direction.

\subsection{Kernel choice and fitting}
\label{sec:kernel-choice}

The kernel defines which scene functions are plausible between samples. We
evaluate an RBF kernel, which assumes very smooth variation, and a
Mat\'ern-$3/2$ kernel, which admits rougher functions and retains more
uncertainty between samples. The formulation is otherwise kernel-agnostic;
anisotropic spatial kernels or material-dependent directional kernels can be
introduced without changing the Bayesian-quadrature equations.

We fit kernel length scales by maximising GP marginal likelihood on
observations sampled from each trained map. This adapts the covariance to the
units and characteristic scale of each scene. The fitted kernel determines
the meaning and absolute scale of the reported variance, so uncertainty
values are compared only under a common kernel model.

\section{Using the Uncertainty}

\subsection{Uncertainty maps}

We form uncertainty images by associating each image region with its
contributing 3-D neighbourhood and evaluating Eq.~\eqref{eq:variance} under
the joint position--direction rendering-aware kernel, so the result
highlights both weak geometric support and surfaces requested from
unfamiliar viewpoints.

\begin{figure*}[t]
    \centering
    \includegraphics[width=\textwidth]{figures/scene_gallery.png}
    \caption{Seven NeRF-Synthetic scenes, one held-out view each, at two
    splat budgets (500 and this project's standard 300{,}000-splat recipe).
    Left to right: ground truth, gsplat reconstruction, $|$error$|$, and raw
    rendering-aware posterior variance (shared log-scale color range across
    every figure in this paper).}
    \label{fig:qualitative}
\end{figure*}

\section{Experiments}
\label{sec:experiments}

We test the two effects that \method{} is designed to measure: incomplete
spatial coverage, from the number of primitives available to represent a
scene, and incomplete directional coverage, from missing viewing angles.

\subsection{Spatial coverage}

We train the lego scene at six splat budgets -- $500$; $10{,}000$;
$30{,}000$; $100{,}000$; $300{,}000$ (this project's standard recipe); and
$1{,}000{,}000$ -- and render the same held-out view against each checkpoint.
Held-out PSNR rises from $21.4\,$dB at $500$ splats to $30.5$, $34.9$,
$38.7$, $40.4$, and $41.0\,$dB at $1{,}000{,}000$, with diminishing returns
past roughly $100{,}000$ splats. Posterior variance tracks the same curve
spatially rather than just in aggregate (Figure~\ref{fig:spatial}): at $500$
splats it is elevated across nearly the whole object, and by
$300{,}000$--$1{,}000{,}000$ splats it concentrates onto a shrinking set of
genuinely under-resolved regions instead of shrinking uniformly.

\begin{figure*}[t]
    \centering
    \includegraphics[width=0.97\textwidth]{figures/lego_splat_sweep.png}
    \caption{Lego, same held-out view, increasing splat budget top to bottom
    ($500$ to $1{,}000{,}000$). Left to right: ground truth, gsplat
    reconstruction, $|$error$|$, and raw rendering-aware posterior variance
    (shared log-scale color range across every figure in this paper).
    Reconstruction quality and BQ uncertainty improve together and saturate
    together.}
    \label{fig:spatial}
\end{figure*}

\subsection{Directional coverage}

We construct a real coverage-gap sweep on the lego scene: symmetric angular
windows of increasing half-width ($0^\circ$ to $75^\circ$) are removed from
the 100 training views, and the same held-out view nearest the gap center is
rendered against each condition's own checkpoint. Reconstruction quality
tracks the missing coverage ($35.1$, $29.2$, $22.9$, $15.1$, and
$16.1\,$dB), and the raw rendering-aware posterior variance grows
monotonically with the gap, by roughly two orders of magnitude: mean
variance $0.82 \to 8.54 \to 22.12 \to 59.45 \to 78.85$ (Figure~\ref{fig:coverage}).

\begin{figure*}[t]
    \centering
    \includegraphics[width=\textwidth]{figures/coverage_uncertainty_sweep.png}
    \caption{Lego, same held-out view, decreasing training-view coverage
    left to right ($0^\circ$ to $75^\circ$ gap half-width). Top to bottom:
    ground truth, gsplat reconstruction, $|$error$|$, and raw rendering-aware
    posterior variance (shared log-scale color range across every figure in
    this paper); the bottom row summarizes the variance's mean and median
    (log scale) over the whole view against gap half-width. Reconstruction
    degrades and BQ uncertainty grows together in the missing-coverage
    region.}
    \label{fig:coverage}
\end{figure*}

\section{Discussion}



\section{Conclusion}



\appendix

Two further experiments, run after the headline spatial- and directional-
coverage results above, on questions this project has not yet resolved
into a paper-level claim: whether the BQ posterior can inform training
itself, not just diagnose a finished checkpoint (\S\ref{app:likelihood}),
and whether kernel choice matters for this method's two central real-data
properties (\S\ref{app:kernels}). Both are reported here for completeness,
at the same standard of honesty as the main results, rather than folded
into the main claims.

\subsection{Training directly under the likelihood}
\label{app:likelihood}

Every result in Section~\ref{sec:experiments} computes BQ posterior
variance \emph{after} training, off a checkpoint fit by ordinary
photometric loss and view-space-gradient-triggered densification. We test
whether the posterior can instead inform training itself, via two
mechanisms already exposed by our trainer: (i) swapping the densification
trigger from gsplat's view-space positional gradient to closed-form BQ
position-only variance, queried at each splat's own position, optionally
floored by opacity so near-transparent, effectively empty splats cannot
outscore splats near real, under-resolved geometry; and (ii) an
uncertainty-weighted Gaussian-NLL auxiliary loss term evaluated on a sparse
grid of real ray-surface points.

We train five matched variants on a real 12-view lego training pool (same
seed and hyperparameters throughout: 2{,}000 initial splats, an
8{,}000-splat budget, 3{,}000 iterations, SSIM-weighted photometric loss,
BQ hyperparameters fit at lego's scale), and evaluate each on both its
training views and a disjoint 15-view subsample of the official held-out
test split. Table~\ref{tab:likelihood} reports the result.

\begin{table}[t]
\centering
\caption{Training directly under the BQ likelihood: matched variants on a
real 12-view lego checkpoint, evaluated on training and held-out views.}
\label{tab:likelihood}
\small
\begin{tabular}{lrrr}
\toprule
Variant & Splats & Train PSNR & Held-out PSNR \\
\midrule
Baseline (gradient densify)          & 3{,}108 & 26.65\,dB & 17.48\,dB \\
BQ-variance densify                  & 3{,}333 & 27.48\,dB & 17.78\,dB \\
NLL loss (gradient densify)          & 3{,}095 & 26.51\,dB & 17.27\,dB \\
BQ-variance densify + NLL            & 3{,}339 & 27.51\,dB & 17.77\,dB \\
BQ-variance densify + opacity floor  & 1{,}181 & 24.94\,dB & 17.64\,dB \\
\bottomrule
\end{tabular}
\end{table}

BQ-variance densification is a real, if modest, held-out win over
gradient-triggered densification ($+0.29\,$dB), consistent whether or not
the NLL term is also active. The NLL auxiliary loss term alone is a
no-op-to-mild-negative ($-0.21\,$dB held-out) at this loss weight and grid
resolution/cadence, plausibly because its variance term does not
differentiate through the BQ posterior itself in this first installment.
Flooring the BQ-variance densification score by opacity is a genuine
efficiency lever rather than a quality tax: $62\%$ fewer splats than
baseline at no held-out PSNR cost. We read this as encouraging but
preliminary: a single scene and iteration budget, not yet a general
training-time claim.

\subsection{Alternative kernels}
\label{app:kernels}

\S\ref{sec:kernel-choice} keeps the position kernel pluggable; the main
results use RBF. We add a third family, a rational quadratic kernel
$k(r) = \left(1 + r^2/(2\alpha l^2)\right)^{-\alpha}$ with $\alpha$ fixed
at $1$ (a continuous scale mixture of RBF bandwidths, plausible for scenes
whose splat density -- and hence natural local length scale -- varies
across the object) and length scale $l$ as its one free parameter, fit by
the same marginal-likelihood procedure as RBF's $\sigma$ and Mat\'ern's
$\rho$. We compare RBF, Mat\'ern-3/2, and rational quadratic, each fit
fresh against real local windows, on all $7$ NeRF-Synthetic scenes used
elsewhere in this paper (chair, drums, ficus, hotdog, lego, mic, ship;
materials is excluded throughout this project for an unrelated,
honestly-reported reason -- its reconstructions stay visibly hazy under
our vanilla-3DGS recipe even at their best held-out view), each at two very
different splat densities (a $300$k-splat checkpoint and a $500$-splat
one), on two properties: whether posterior variance tracks real local
splat sparsity (distance to the $8$th-nearest real neighbor, at $150$ real
splat positions per checkpoint), and whether it is calibrated against real
held-out rendering error (up to $480$ depth-unprojected query points per
checkpoint from held-out views, scored by correlation and by a
Gaussian-NLL-style proper scoring rule, lower is better). The per-checkpoint
local-window radius (tuned once, on lego, in an earlier single-scene
version of this experiment) was re-checked directly against real neighbor
counts on all $7$ scenes before reuse: it transfers without retuning
everywhere (every scene/checkpoint combination keeps $\geq 6$ real
neighbors for at least $86\%$ of sampled windows).

\begin{table*}[t]
\centering
\caption{Kernel-family comparison at fitted bandwidths, on a dense
($300$k-splat) and a sparse ($500$-splat) checkpoint of the chair scene.
Sparsity $r$ is between posterior variance and distance to the
$8$th-nearest real neighbor (a working signal is positive); ``ratio''
normalizes variance by its own prior (amplitude-independent) variance.
Calibration $r$/mean NLL are against real held-out rendering error. The
same columns are repeated for the other $6$ scenes in
Tables~\ref{tab:kernels-drums}--\ref{tab:kernels-ship}.}
\label{tab:kernels}
\small
\begin{tabular}{llrrrrr}
\toprule
Checkpoint & Kernel & Fitted param.\ & Sparsity $r$ (raw) & Sparsity $r$ (ratio) & Calib.\ $r$ & Mean NLL \\
\midrule
Dense (300k) & RBF                & $\sigma=0.0608$ & $+0.473$ & $+0.508$          & $+0.120$ & $\mathbf{21.9}$ \\
Dense (300k) & Mat\'ern-3/2       & $\rho=0.0180$   & $+0.693$ & $\mathbf{+0.779}$ & $+0.050$ & $7635.6$ \\
Dense (300k) & Rational quadratic & $l=0.0100$      & $+0.706$ & $+0.768$          & $+0.041$ & $8808.3$ \\
Sparse (500) & RBF                & $\sigma=0.0952$ & $+0.647$ & $+0.785$          & $-0.077$ & $\mathbf{-0.29}$ \\
Sparse (500) & Mat\'ern-3/2       & $\rho=0.0155$   & $+0.320$ & $+0.879$          & $-0.146$ & $1246.9$ \\
Sparse (500) & Rational quadratic & $l=0.0120$      & $+0.342$ & $\mathbf{+0.885}$ & $-0.149$ & $470.0$ \\
\bottomrule
\end{tabular}
\end{table*}

\begin{table*}[t]
\centering
\caption{Kernel-family comparison, drums scene. Columns as in Table~\ref{tab:kernels}.}
\label{tab:kernels-drums}
\small
\begin{tabular}{llrrrrr}
\toprule
Checkpoint & Kernel & Fitted param.\ & Sparsity $r$ (raw) & Sparsity $r$ (ratio) & Calib.\ $r$ & Mean NLL \\
\midrule
Dense (300k) & RBF                & $\sigma=0.0546$ & $-0.062$ & $-0.213$          & $-0.051$ & $\mathbf{22.5}$ \\
Dense (300k) & Mat\'ern-3/2       & $\rho=0.0109$   & $-0.076$ & $+0.183$          & $-0.077$ & $2.93\times 10^{4}$ \\
Dense (300k) & Rational quadratic & $l=0.0060$      & $-0.082$ & $\mathbf{+0.245}$ & $-0.084$ & $3.40\times 10^{4}$ \\
Sparse (500) & RBF                & $\sigma=0.1111$ & $+0.718$ & $+0.639$          & $+0.059$ & $\mathbf{-0.10}$ \\
Sparse (500) & Mat\'ern-3/2       & $\rho=0.0623$   & $+0.642$ & $+0.707$          & $-0.120$ & $25.0$ \\
Sparse (500) & Rational quadratic & $l=0.0370$      & $+0.639$ & $\mathbf{+0.725}$ & $-0.133$ & $25.1$ \\
\bottomrule
\end{tabular}
\end{table*}

\begin{table*}[t]
\centering
\caption{Kernel-family comparison, ficus scene. Columns as in Table~\ref{tab:kernels}.}
\label{tab:kernels-ficus}
\small
\begin{tabular}{llrrrrr}
\toprule
Checkpoint & Kernel & Fitted param.\ & Sparsity $r$ (raw) & Sparsity $r$ (ratio) & Calib.\ $r$ & Mean NLL \\
\midrule
Dense (300k) & RBF                & $\sigma=0.0578$ & $-0.193$ & $-0.308$          & $+0.112$ & $\mathbf{60.0}$ \\
Dense (300k) & Mat\'ern-3/2       & $\rho=0.0131$   & $+0.185$ & $-0.269$          & $-0.026$ & $4.47\times 10^{4}$ \\
Dense (300k) & Rational quadratic & $l=0.0077$      & $+0.114$ & $\mathbf{-0.193}$ & $+0.047$ & $4.32\times 10^{4}$ \\
Sparse (500) & RBF                & $\sigma=0.1033$ & $+0.769$ & $+0.751$          & $+0.149$ & $\mathbf{0.33}$ \\
Sparse (500) & Mat\'ern-3/2       & $\rho=0.0071$   & $+0.724$ & $+0.884$          & $+0.151$ & $2.45\times 10^{4}$ \\
Sparse (500) & Rational quadratic & $l=0.0050$      & $+0.725$ & $\mathbf{+0.893}$ & $+0.150$ & $1.11\times 10^{4}$ \\
\bottomrule
\end{tabular}
\end{table*}

\begin{table*}[t]
\centering
\caption{Kernel-family comparison, hotdog scene. Columns as in Table~\ref{tab:kernels}.}
\label{tab:kernels-hotdog}
\small
\begin{tabular}{llrrrrr}
\toprule
Checkpoint & Kernel & Fitted param.\ & Sparsity $r$ (raw) & Sparsity $r$ (ratio) & Calib.\ $r$ & Mean NLL \\
\midrule
Dense (300k) & RBF                & $\sigma=0.1110$ & $-0.341$ & $\mathbf{+0.582}$ & $+0.202$ & $\mathbf{2298.8}$ \\
Dense (300k) & Mat\'ern-3/2       & $\rho=0.0455$   & $-0.420$ & $+0.503$          & $+0.122$ & $7191.2$ \\
Dense (300k) & Rational quadratic & $l=0.0180$      & $-0.425$ & $+0.516$          & $+0.113$ & $1.06\times 10^{4}$ \\
Sparse (500) & RBF                & $\sigma=0.1204$ & $+0.557$ & $+0.575$          & $+0.123$ & $\mathbf{-0.24}$ \\
Sparse (500) & Mat\'ern-3/2       & $\rho=0.0531$   & $+0.497$ & $+0.867$          & $-0.005$ & $62.8$ \\
Sparse (500) & Rational quadratic & $l=0.0317$      & $+0.494$ & $\mathbf{+0.877}$ & $-0.008$ & $61.2$ \\
\bottomrule
\end{tabular}
\end{table*}

\begin{table*}[t]
\centering
\caption{Kernel-family comparison, lego scene. Columns as in Table~\ref{tab:kernels}. (This scene's checkpoints are the same ones used throughout the main paper's figures.)}
\label{tab:kernels-lego}
\small
\begin{tabular}{llrrrrr}
\toprule
Checkpoint & Kernel & Fitted param.\ & Sparsity $r$ (raw) & Sparsity $r$ (ratio) & Calib.\ $r$ & Mean NLL \\
\midrule
Dense (300k)  & RBF                & $\sigma=0.0703$ & $-0.306$ & $-0.291$          & $+0.120$ & $\mathbf{9.80}$ \\
Dense (300k)  & Mat\'ern-3/2       & $\rho=0.0186$   & $-0.441$ & $-0.019$          & $+0.062$ & $1878.0$ \\
Dense (300k)  & Rational quadratic & $l=0.0067$      & $-0.458$ & $\mathbf{+0.191}$ & $+0.064$ & $6062.0$ \\
Sparse (500)  & RBF                & $\sigma=0.0988$ & $+0.766$ & $+0.811$          & $-0.055$ & $\mathbf{0.08}$ \\
Sparse (500)  & Mat\'ern-3/2       & $\rho=0.0406$   & $+0.783$ & $+0.916$          & $-0.112$ & $87.8$ \\
Sparse (500)  & Rational quadratic & $l=0.0238$      & $+0.783$ & $\mathbf{+0.921}$ & $-0.112$ & $85.4$ \\
\bottomrule
\end{tabular}
\end{table*}

\begin{table*}[t]
\centering
\caption{Kernel-family comparison, mic scene. Columns as in Table~\ref{tab:kernels}.}
\label{tab:kernels-mic}
\small
\begin{tabular}{llrrrrr}
\toprule
Checkpoint & Kernel & Fitted param.\ & Sparsity $r$ (raw) & Sparsity $r$ (ratio) & Calib.\ $r$ & Mean NLL \\
\midrule
Dense (300k) & RBF                & $\sigma=0.0588$ & $-0.368$ & $-0.294$          & $+0.141$ & $\mathbf{312.1}$ \\
Dense (300k) & Mat\'ern-3/2       & $\rho=0.0110$   & $-0.443$ & $\mathbf{+0.046}$ & $+0.013$ & $3.69\times 10^{6}$ \\
Dense (300k) & Rational quadratic & $l=0.0052$      & $-0.462$ & $-0.083$          & $+0.016$ & $6.11\times 10^{6}$ \\
Sparse (500) & RBF                & $\sigma=0.0997$ & $+0.215$ & $+0.034$          & $+0.127$ & $\mathbf{1.31\times 10^{8}}$ \\
Sparse (500) & Mat\'ern-3/2       & $\rho=0.0506$   & $+0.228$ & $+0.139$          & $-0.019$ & $9.28\times 10^{8}$ \\
Sparse (500) & Rational quadratic & $l=0.0292$      & $+0.231$ & $\mathbf{+0.172}$ & $-0.030$ & $9.28\times 10^{8}$ \\
\bottomrule
\end{tabular}
\end{table*}

\begin{table*}[t]
\centering
\caption{Kernel-family comparison, ship scene. Columns as in Table~\ref{tab:kernels}.}
\label{tab:kernels-ship}
\small
\begin{tabular}{llrrrrr}
\toprule
Checkpoint & Kernel & Fitted param.\ & Sparsity $r$ (raw) & Sparsity $r$ (ratio) & Calib.\ $r$ & Mean NLL \\
\midrule
Dense (300k) & RBF                & $\sigma=0.0683$ & $\mathbf{+0.824}$ & $\mathbf{+0.750}$ & $+0.220$ & $\mathbf{13.7}$ \\
Dense (300k) & Mat\'ern-3/2       & $\rho=0.0193$   & $+0.427$          & $+0.609$          & $+0.037$ & $3281.8$ \\
Dense (300k) & Rational quadratic & $l=0.0095$      & $+0.432$          & $+0.637$          & $+0.005$ & $5696.4$ \\
Sparse (500) & RBF                & $\sigma=0.1346$ & $\mathbf{+0.760}$ & $\mathbf{+0.927}$ & $-0.012$ & $\mathbf{-0.43}$ \\
Sparse (500) & Mat\'ern-3/2       & $\rho=0.0698$   & $+0.384$          & $+0.882$          & $-0.020$ & $3.92$ \\
Sparse (500) & Rational quadratic & $l=0.0441$      & $+0.382$          & $+0.887$          & $-0.019$ & $3.50$ \\
\bottomrule
\end{tabular}
\end{table*}

No single family dominates every metric, checkpoint, or scene
(Tables~\ref{tab:kernels}--\ref{tab:kernels-ship}), but one part of the
lego-only finding turns out to be fully universal and another turns out to
be scene-dependent in an informative way.

RBF is the best-calibrated family by the Gaussian-NLL score in
\emph{every one} of the $14$ scene/checkpoint combinations, typically by
two or more orders of magnitude and up to four in the worst case (mic,
where all three families are in fact badly miscalibrated in absolute
terms -- real held-out error there is large relative to any family's
fitted variance -- but RBF is still the least catastrophic by roughly
$1000\times$). Mat\'ern/rational quadratic's marginal-likelihood fits
consistently land on tighter bandwidths that explain local color
smoothness well but produce overconfident (too-small) variance at real
held-out points, punished heavily by the $1/\text{var}$ term -- the same
mechanism identified on lego alone, now confirmed project-wide. The
calibration \emph{correlation} itself stays weak and sign-inconsistent
across every family and every scene, exactly as on lego: none of the three
families should be read as genuinely tracking held-out error in the
correlation sense.

The sparsity-tracking story is less clean than lego alone suggested. On
the sparse ($500$-splat) checkpoint, raw variance has the expected
positive sign for every family on all $7$ scenes ($21/21$), and the
amplitude-normalized ratio is won by rational quadratic on $6$ of $7$
scenes (all but ship, where RBF is narrowly ahead, $+0.927$ vs.\ $+0.887$)
-- a broad, if not perfect, replication of the lego finding. On the dense
($300$k-splat) checkpoint, however, lego's ``wrong sign for every family''
result does \emph{not} generalize: raw variance keeps the intended
positive sign on chair and ship, is split by family on ficus, and is
wrong-signed for every family only on drums, hotdog, lego, and mic. This
tracks a real per-scene difference we checked directly: the number of
real neighbors inside a fixed-radius query window at this checkpoint's
density ranges from a median of $\sim\!21$ (ficus) to $\sim\!9{,}000$
(mic) even though every scene shares the same training recipe and window
radius -- it is specifically the high-neighbor-count scenes that
saturate the engine's neighbor cap and let local opacity confound raw
variance, not a universal property of dense checkpoints. Reading via the
ratio on the dense checkpoint accordingly gives a more mixed picture too:
rational quadratic is best on $3$ of $7$ scenes (drums, ficus, lego),
Mat\'ern on $2$ (chair, mic), and RBF on $2$ (hotdog, ship). We read this
as the same genuine trade-off as before, now on firmer footing for
calibration and a real but noisier one for sparsity-tracking: RBF where a
well-behaved, non-catastrophic variance estimate matters most (e.g.\ any
NLL-style loss, including \S\ref{app:likelihood}'s), rational quadratic as
the best single default where sparsity/coverage-tracking is the goal,
particularly on genuinely sparse checkpoints -- but not a law that holds
irrespective of scene content at high splat density.

\bibliographystyle{IEEEtran}
\bibliography{references}

\end{document}